\chapter{Literatuurstudie}
\label{ch:literatuur}

% Schrijf per topic, niet per paper. Elke sectie volgt hetzelfde patroon:
%   1. Wat is er bekend? (synthese over meerdere bronnen heen)
%   2. Hoe wordt het gemeten / met welke data en methoden?
%   3. Wat zijn de beperkingen?
%   4. Gap: wat ontbreekt en welke onderzoeksvraag pakt dat aan?
% Topiccodes (A-F) komen overeen met het tabblad 'Topics' in
% Literatuur_studie/bronnen_literatuurstudie.xlsx.

% Inleiding: doel van het hoofdstuk en opbouw (welke secties, waarom die volgorde).

\section{Fysieke prestatie-indicatoren in het voetbal}
\label{sec:lit-indicatoren}
% Topic A -> OV1
% Totale afstand, high-speed running, sprints, topsnelheid, acceleraties/deceleraties,
% richtingsveranderingen, luchtduels. Verschillen per positie en niveau.

\section{Meten van fysieke prestaties}
\label{sec:lit-meten}
% Topic B -> OV1 (stabiliteit)
% Trackingsystemen (optisch, GPS/LPS, broadcast tracking zoals SkillCorner), validiteit,
% variabiliteit van wedstrijd tot wedstrijd, normalisatie (per 90 min, per positie,
% balbezit), drempelwaarden voor snelheidszones.
%
% Synchronisatie en preprocessing van event- en trackingdata:
% - \citet{oonk2026databallpy}: open-source Python-pakket (DataBallPy) dat event- en
%   trackingdata in één Game-object bundelt en synchroniseert met een Needleman-Wunsch-
%   algoritme. Naïeve timestamp-alignering wijkt 1.82 (±4.06) s af; correcte synchronisatie
%   verlaagt de Brier loss van een xG-model van 0.096 naar 0.082. Relevant omdat fysieke
%   metrieken per balbezit/fase (normalisatie, Taak 3) en de koppeling mismatch <-> kans (OV3)
%   een correcte synchronisatie vereisen. Berekent ook afstand in snelheids-/acceleratiezones.

\section{Spelersprofielen en rolclustering}
\label{sec:lit-profielen}
% Topic C -> OV1
% Regel- en percentielgebaseerde profielen versus clustering (k-means, GMM, ...),
% keuze van het aantal profielen, interpreteerbaarheid, stabiliteit.

\section{Matchups en mismatches}
\label{sec:lit-mismatches}
% Topic D -> OV2
% Wie staat tegenover wie (marking assignment), duelanalyse, ruimtecontrole (pitch control),
% vergelijking op niveau van duel, linie en team.
%
% - Implementaties van pitch control (Fernandez & Bornn, 2018) en pressure zijn beschikbaar in
%   \citet{oonk2026databallpy}; mogelijke bouwstenen voor mismatches op duel- en linieniveau.

\section{Koppeling met wedstrijdgebeurtenissen}
\label{sec:lit-uitkomst}
% Topic E -> OV3
% Expected goals, waardering van acties, verband tussen fysieke output en uitkomst.
%
% - Keuze van de uitkomstmaat: \citet{bergkamp2019methodological} pleiten voor individuele,
%   positiespecifieke prestatiematen (gewogen events, positiedata) i.p.v. prestatieniveau,
%   en wijzen erop dat individuele prestatie afhangt van ploegmaats en tegenstanders.

\section{Jeugdvoetbal, talentdetectie en ethiek}
\label{sec:lit-jeugd}
% Topic F -> context en ethische reflectie
% Maturatie, relative age effect, talentidentificatie, gebruik van data van minderjarigen.
%
% Expertoordeel vs. actuariële combinatie van informatie:
% - \citet{peringa2025uniform}: Lens Model. Scouts combineren cues inconsistent en
%   suboptimaal; het lineaire model van een scout is accurater dan de scout zelf
%   (ra = .62, rm = .68, Re = .77). Argument voor consistente, gewogen profielscores
%   i.p.v. een holistisch oordeel. Ook bruikbaar als opzet voor de expertvalidatie
%   met RBFA-analisten (Taak 7).
%
% Relative age effect in de Belgische nationale ploegen:
% - \citet{billen2022careers} (bachelorproef KU Leuven, niet peer-reviewed; RBFA-data, 29.631
%   datapunten): significant RAE in U15-U18, niet in U19, U21 en de Rode Duivels; geen verschil
%   per positie. Selectie bij de U21 (OR = 4.55) en de Olympische ploeg (OR = 6.89) voorspelt
%   een selectie bij de Rode Duivels. Relevant: bij jeugdspelers kunnen maturiteit en
%   geboortekwartaal fysieke profielen vertekenen (RBFA-project 'Futures'). Voorzichtig citeren.
%
% Methodologie van talentidentificatie:
% - \citet{bergkamp2019methodological}: vier methodologische problemen vanuit de
%   selectiepsychologie: (1) het criterium is meestal het prestatieniveau of de selectie zelf
%   i.p.v. individuele wedstrijdprestatie; (2) geïsoleerde testen zijn low-fidelity 'signs',
%   wedstrijdgedrag (event-/positiedata) zijn high-fidelity 'samples' en voorspellen beter;
%   (3) range restriction in voorgeselecteerde groepen onderschat verbanden; (4) base rate
%   en selectieratio bepalen de utiliteit (voorbeeld: r = .4 -> success ratio 5.3%).
%   Relevant: fysieke profielen uit trackingdata = samples; bij RBFA-spelers (sterk
%   voorgeselecteerd) zijn fysieke verschillen klein -> populatie expliciet vermelden;
%   uitkomstmaat voor OV3 individueel en positiespecifiek kiezen en eerst valideren.
% - \citet{bareli2024feudal} (enkel abstract): talentselectie combineert het subjectieve
%   'coach's eye' en objectieve testformules; de auteurs pleiten voor big data / sport
%   analytics als derde aanpak die beide verbindt. Hoogstens één zin context.
% - \citet{barraclough2022methodological} (enkel abstract): narrative review over TID in
%   teamsporten; pleit voor longitudinale, multidisciplinaire en ecologisch valide designs.
%   Bevestigt \citet{bergkamp2019methodological}: fysieke profielen uit wedstrijden, opgevolgd
%   over wedstrijden en seizoenen (OV1-stabiliteit).
% - \citet{bennett2026talent}: real-time Delphi met 20 experts over 15 aanbevelingen voor TID/TD.
%   90% vindt dat selectors hun intuïtie moeten aanvullen met gevalideerde objectieve data, maar
%   die data moet betekenisvol en contextueel zijn (prestatiedata bevat 'te veel ruis'). Indicatoren
%   longitudinaal opvolgen: 75% eens, maar slechts 10% haalbaar. Argument voor interpreteerbare
%   profielen en een automatische pijplijn; Delphi-opzet (consensus = IQR <= 1.25) als optie
%   voor de expertvalidatie met RBFA-analisten (Taak 7).

\section{Synthese}
\label{sec:lit-synthese}
% Tabel met technieken: sterktes, beperkingen, benodigde data (= verifieerbaar resultaat Taak 1).
% Welke technieken neem je mee naar het baselinemodel (Taak 3) en waarom?
% Herformuleer de gaps en koppel ze expliciet aan OV1-OV3.
% Methodologische aandachtspunten uit \citet{houtmeyers2021managing}: datareductie bij veel
% gecorreleerde metrieken (bv. PCA), onzekerheid rapporteren, correlatie is geen causaliteit
% (OV3), in-sample vs. out-of-sample evaluatie tegen overfitting.

% Voorbeeld van een overzichtstabel:
% \begin{table}[htbp]
%   \centering
%   \caption{Overzicht van technieken uit de literatuur.}
%   \label{tab:lit-technieken}
%   \begin{tabular}{@{}llll@{}}
%     \toprule
%     Techniek & Sterktes & Beperkingen & Benodigde data \\
%     \midrule
%     ... & ... & ... & ... \\
%     \bottomrule
%   \end{tabular}
% \end{table}
